\documentclass[12pt, a4paper]{article}

% ============================================================
%  Preamble (mirrors leak-whitepaper.tex so the two documents
%  share one toolchain: latexmk -> pdflatex + biber)
% ============================================================
\usepackage[utf8]{inputenc}
\usepackage[T1]{fontenc}
\usepackage{lmodern}
\usepackage{microtype}
\usepackage{amsmath, amssymb, amsthm}
\usepackage{graphicx}
\usepackage{booktabs}
\usepackage{array}
\usepackage{longtable}
\usepackage{listings}
\usepackage{xcolor}
\usepackage{enumitem}
\usepackage{titlesec}
\usepackage{abstract}
\usepackage{caption}
\usepackage{float}
\usepackage[margin=2.8cm]{geometry}
\usepackage[backend=biber, style=numeric, urldate=long]{biblatex}
\usepackage{hyperref} % load last

\addbibresource{sample.bib}
\addbibresource{benchmark-report.bib}

\hypersetup{
    colorlinks=true,
    linkcolor=blue!60!black,
    urlcolor=blue!60!black,
    citecolor=green!50!black
}

\emergencystretch=3em % long \texttt run/theorem names don't hyphenate
\titleformat{\section}{\large\bfseries}{\thesection.}{0.6em}{}
\titleformat{\subsection}{\normalsize\bfseries}{\thesubsection}{0.6em}{}

\definecolor{codebg}{HTML}{F5F5F5}
\definecolor{codekeyword}{HTML}{2E4DA0}
\lstdefinelanguage{Lean4}{
  keywords={theorem, lemma, def, let, have, intro, cases, induction, exact,
            apply, refine, rw, simp, ring, linarith, nlinarith, omega, by,
            sorry, import, open, Type, Prop},
  keywordstyle=\color{codekeyword}\bfseries,
  sensitive=true,
  comment=[l]{--},
}
\lstset{language=Lean4, backgroundcolor=\color{codebg}, basicstyle=\ttfamily\small,
        breaklines=true, frame=single, framerule=0pt, xleftmargin=6pt,
        literate={∃}{{$\exists$}}1 {∧}{{$\wedge$}}1 {→}{{$\to$}}1
                 {⊗}{{$\otimes$}}1 {⊕}{{$\oplus$}}1 {∀}{{$\forall$}}1}

% Shorthand for run labels / ids / commits
\newcommand{\run}[1]{\texttt{#1}}
\newcommand{\cmark}{$\checkmark$}

\title{\textbf{To what extent does scaffolding improve proving capacity of automated theorem
provers driven by large agentic systems?}}
\author{Mikael Bashir - 02382080}
\date{14 August 2026}

\begin{document}
\maketitle

\begin{abstract}
\noindent
We report results from three benchmark runs of the Leak \cite{leakabout} automated
theorem-proving stack on FATE-X \cite{fatex2025}, a 100-problem graduate-level
abstract and commutative-algebra benchmark in Lean~4 \cite{demoura2021lean4}.The headline 
finding is negative for our central design hypothesis: on identical time budgets and the same base agent, a single
flat agent with compiler and search access (Control~II) solved
\textbf{38/98 (38.8\%)}, outperforming the best decomposition pipeline
Ultra Fleeting, \textbf{28/98, (28.6\%)}. Decomposition contributed only two problems
the flat agent could not solve. We analyse where each approach wins {put an in paper link here}, quantify a
difficulty cliff in FATE-X's top quartile (2/25 for both arms) {put an in paper link here}, document the
operational incidents that shaped the design, including the MCP services we made 
extensive use of {put an in paper link here}, and lay out the follow-up experiments the results motivate or 
that could not be completed due to budget {put an in paper link here}.
\end{abstract}

\tableofcontents
\newpage

% ============================================================
\section{Introduction}
% ============================================================

There has been extensive work on LLMs, specifically on Automated Theorem Proving systems 
driven by LLMs, with a variety of architecture designs, both in the development of the LLM itself,
and the orchestration of the LLM. We are also seeing the number of agentic systems being made publically available
increase rapidly, which raises the natural question "Do the same techniques for ATP on LLMs apply to Agentic systems, and to what extent". 
We present the Leak stack, an agentic theorem-proving system for Lean~4 + Mathlib
\cite{demoura2021lean4,mathlib2020} that tries to answer exactly this, armed with select, specialist remote verification services
(``Leak~I--XIV'') and a local bridge that drives LLM agents against them
\cite{competemath,hfspaces}. Its design bet, shared with a line of recent work
on hierarchical and decomposition-based proving
\cite{lample2022hypertree,deepseekproverv2,decomposition_openreview,davis2025goedelspoetry},
is that hard theorems are best attacked by splitting them into sub-goals proved
by parallel sub-agents. This report evaluates that bet empirically, while also 
testing if this still holds in large agentic systems such as Claude, which already have secret
scaffolding built into them {put the sdk harness ref here, as well as the agent terminology definitions ref}.

We ran several full benchmark passes on a sanitised {put in paper link to how we sanitised here} FATE-X
\cite{fatex2025,benchmarklist_fatex}, recording every attempt as a structured
research row. The experiments were designed to answer three questions:

\begin{enumerate}[itemsep=1pt]
  \item \textbf{Absolute proving capacity.} What fraction of an algebra
        benchmark, exceeding PhD qualifying exam difficulty,
        can Leak close under fixed time budgets (FATE-X)?
  \item \textbf{Adding scaffolding vs.\ Relying on built-in scaffolding.} Does the lifted architecture
        (blueprint\,$\to$\,node-provers\,$\to$\,refinement) beat a
        single continuous agent given the same model, budget, and select proving tools?
  \item \textbf{Does Sonnet 5 benefit significantly when given pantograph?} While we armed many runs
        with an improved pantograph service {put link here to that service's repo, at Leak II}, 
        we always took for granted its benefit - but in the context of agentic systems, how much
        does pantograph actually buy.
\end{enumerate}

Questions 1 is objectively answered below, and we provide evidences below for questions 2 and 3.

% ============================================================
\section{Related work}
% ============================================================

Whole-proof and agentic provers are typically evaluated on miniF2F
\cite{zheng2022minif2f}, PutnamBench \cite{tsoukalas2024putnambench}, or
competition sets targeted by systems such as GPT-f \cite{polu2020gptf},
HyperTree Proof Search \cite{lample2022hypertree}, AlphaProof
\cite{alphaproof2024}, DeepSeek-Prover-V2 \cite{deepseekproverv2}, and the
G\"odel-Prover line \cite{davis2025goedelspoetry}. FATE-X \cite{fatex2025}
is deliberately harder and narrower: 100 abstract/commutative-algebra problems
drawn from graduate textbooks, qualifying exams, and research literature,
ordered by increasing difficulty. Its authors report frontier systems solving
only single-digit counts on the upper half \cite{benchmarklist_fatex}, which
matches the cliff we observe (\S\ref{sec:difficulty}). 

Decomposition-style pipelines --- sketching a proof outline and delegating
sub-goals --- descend from Draft--Sketch--Prove-like designs and are the core of
DeepSeek-Prover-V2's subgoal decomposition \cite{deepseekproverv2} and related
hierarchical searches \cite{lample2022hypertree,decomposition_openreview}. We test an already
researched and documented pipeline, Goedel-Architect \cite{goedelarchitect}, changing the driving
backbone --- instead of a fast and small LLM, such as the DeepSeek-V4-Flash backbone used by
Goedel-Architect, we drive with a large agent (Claude Sonnet~5 Agent \cite{anthropic2026sonnet5}),
referring to this design as Ultra Fleeting through out the paper;
as far as we are aware the team behind Goedel Architect have not made available the critical
scaffolding on DeepSeek-V4-Flash, so we try to replicate as closely as possible from their report.
The main features are a blueprint stage that compiles a skeleton of \texttt{sorry}-holes, 
node provers that fill holes as ReAct-style tool agents \cite{yao2023react} against live compiler service feedback, 
and a refinement stage restructures the skeleton between iterations. We also experiment with 
Pantograph \cite{aniva2024pantograph,pypantograph}, extended with a
snapshot-based ``ghost state'' model \cite{shen2026snapshot,leakii2026upgrade}
that lets many agents manage proof states on a single Lean daemon, without expensive re-elaboration. 
This service was used for Control-II runs and other experimental runs not presented in this report, 
but it wasn't used for the Ultra Fleeting benchmark, in line with the design of
the borrowed architecture. In the further-works section (\S\ref{sec:nextsteps}), we discuss some
interesting experiments that could not be completed before the writing of this report.

% ============================================================
\section{Experimental setup}
% ============================================================

\subsection{Benchmark and normalization}
\label{sec:benchmark}

FATE-X contains 100 Lean files pinning Lean \texttt{v4.28.0} \cite{fatex2025}.
Our ingestion (\run{build-fatex.mjs}) normalizes each item: imports and the
\texttt{namespace ProblemN} wrapper are dropped, doc comments become the
informal statement, \texttt{section} scoping is kept, and the 38 problems with
supporting declarations keep them inside the statement. Problems are enqueued
in the benchmark's own difficulty order, so the item index doubles as a
difficulty rank. Two statements (\run{fatex\_009}, \run{fatex\_041}) do not
elaborate under our newer toolchains and were excluded from the full passes,
leaving $n=98$.

\subsection{Verification infrastructure}
\label{sec:verifiers}

All scoring is by independent re-verification, never by agent self-report.
Two verifier groups were in play, and they differ by toolchain:

\begin{itemize}[itemsep=1pt]
  \item \textbf{Leak~IV group} (Lean + Mathlib \texttt{v4.29.1}): gates every
        Stronghold-family and control run. Verification is a daemon compile of
        the full submitted script; a proof counts only if it compiles
        \texttt{sorry}-free and proves the exact target signature.
  \item \textbf{Leak~XI/XII/XIV group} (Lean + Mathlib \texttt{v4.32.0}):
        serves the architect pipelines (River, Ultra) --- XII compiles and
        elaborates, XIV verifies full scripts, XI provides Loogle/Moogle-style
        search \cite{loogle,moogle}.
\end{itemize}

\noindent
The two full passes therefore ran on \emph{different} Mathlib versions
(\S\ref{sec:validity} discusses the confound). Search tooling on the
Stronghold/control side is Leak~I \cite{loogle}; all services follow the
MCP protocol \cite{anthropic2024mcp}.

\subsection{Strategy catalogue}
\label{sec:strategies}

\begin{description}[itemsep=2pt]
  \item[River (Grok driver).] The architect pipeline driven over the xAI API
        \cite{xai2025grok41}. The benchmarked variant, \run{river-vintage},
        adds an oversight watcher team (interceptor + mechanic + consultant)
        around the loop.
  \item[Ultra (Claude CLI driver).] The same pipeline driven by the local
        Claude CLI \cite{anthropic2026sonnet5}; \run{ultra-fleeting} is the
        plain arm: no dead-end ledger, no NL seed, fresh context per stage.
  \item[Stronghold (flat-skeleton family).] A planner emits one skeleton of
        \texttt{have}-holes; minions fill holes against the Leak~II ghost-state
        daemon \cite{shen2026snapshot}. Arms: \run{have-tree} (``Dark'',
        isolated minions), \run{have-surround} (continuous parallel work-queue),
        \run{finality-1} (Surround + timed refinement + a proposition-keyed
        lemma bank), \run{stronghold-force}, \run{stronghold-forte}, and
        \run{stronghold-impenetrable} (three-phase clock; enqueued but never
        scored this campaign).
  \item[Controls.] \S\ref{sec:controls}.
\end{description}

\subsection{Protocol}
\label{sec:protocol}

One problem per worker; every attempt is bounded by wall clock only --- no
per-call caps. Budgets grew as the campaign matured: 20\,min for the first
probes (28--29~July), 30\,min mid-campaign, and 60\,min for the two full passes
and all later runs (Table~\ref{tab:ledger}). Architect arms additionally carry
a refinement-iteration budget (5 for \run{river-vintage}, whose per-attempt
cutoff is already halved by the watcher variant; 8 for every other architect
arm).

Failure triage distinguishes infrastructure failures from problem failures:
catastrophic errors (quota, bridge down) pause the pass and \emph{release} the
claim without burning an attempt; problem-specific errors retry once, then mark
the item \emph{skipped}, which excludes it from the denominator rather than
counting as a miss. A watchdog aborts streams wedged 5 minutes past the
bridge's own deadline --- a stuck-socket backstop, not a prover cap. Every
finished attempt files a research row (tagged
\run{benchmark:fatex:<runId>}) into the same tables as production runs,
with full per-attempt metrics.

\subsection{Metrics}

Pass rate is $\text{proved}/\text{scored}$, where scored $=$ proved $+$
unsolved $+$ refuted (skipped excluded). Costs are the CLI-reported
\texttt{total\_cost\_usd} summed per item; for Claude-driver runs this is a
\emph{notional API-equivalent} figure (the operator runs on a flat-rate
subscription), while Grok-driver costs are real API spend. Times are
per-item wall clock from claim to scoring.

% ============================================================
\section{Run ledger}
% ============================================================

Table~\ref{tab:ledger} lists every run in the campaign. Small runs are
deliberate probes (the protocol favours cheap $n{\le}5$ probes before
committing a 98-item pass); their statistics are indicative only.

\begin{table}[H]
\centering\small
\caption{All benchmark runs, 28~July -- 14~August 2026. P/U = proved/unsolved
(scored items); rate = P/(P+U+refuted). Budget is per-item wall clock. All
runs used claude-sonnet-5 except \run{LRV-Indomitable} (opus-4.8, Grok-driven
pipeline). $^{\dagger}$notional API-equivalent cost (\S\ref{sec:cost});
$^{*}$run in progress (95 items pending).}
\label{tab:ledger}
\begin{tabular}{@{}lllrrrr@{}}
\toprule
Run & Strategy & Budget & P & U & Rate & Cost (\$) \\
\midrule
\run{LRV-Indomitable} & \run{river-vintage} & 20m & 0 & 5 & 0\% & 1.08 \\
\run{LUF-Flowing} & \run{ultra-fleeting} & 20m & 3 & 2 & 60\% & 15.16$^{\dagger}$ \\
\run{LSD-Artist} & \run{have-tree} & 20m & 0 & 1 & 0\% & 3.65$^{\dagger}$ \\
\run{LSS-Imminent} & \run{have-surround} & 30m & 3 & 1 & 75\% & 22.30$^{\dagger}$ \\
\run{LFI-Never-Ending} & \run{finality-1} & 30m & 1 & 1 & 50\% & 5.18$^{\dagger}$ \\
\run{LFI-Regime} & \run{finality-1} & 30m & 0 & 1 & 0\% & 4.64$^{\dagger}$ \\
\run{LFI-Ambition} & \run{finality-1} & 30m & 0 & 0 & --- & 5.01$^{\dagger}$ \\
\run{LFI-Hope} & \run{finality-1} & 30m & 0 & 2 & 0\% & 7.17$^{\dagger}$ \\
\run{LFI-All-Seeing} & \run{finality-1} & 30m & 0 & 2 & 0\% & 9.64$^{\dagger}$ \\
\run{LSF-Patient} & \run{stronghold-force} & 60m & 0 & 1 & 0\% & 4.86$^{\dagger}$ \\
\run{LFI-Undone} & \run{stronghold-force} & 60m & 0 & 1 & 0\% & 19.45$^{\dagger}$ \\
\run{LUF-Anonymous} & \run{ultra-fleeting} & 60m & 28 & 70 & 28.6\% & 825.53$^{\dagger}$ \\
\run{LSF-Curiosity} & \run{stronghold-forte} & 60m & 0 & 1 & 0\% & 11.87$^{\dagger}$ \\
\run{LC-Intuition} & \run{control-oneshot} & 30m & 2 & 3 & 40\% & 7.13$^{\dagger}$ \\
\run{LC-II-Irrationality} & \run{control-oneshot-2} & 60m & 38 & 60 & 38.8\% & 878.57$^{\dagger}$ \\
\run{LSI-Unchartered} & \run{stronghold-impenetrable} & 60m & 0 & 0 & --- & --- \\
\run{LC-III-Fair}$^{*}$ & \run{control-oneshot-3} & 60m & 1 & 2 & 33\% & 11.16$^{\dagger}$ \\
\midrule
\multicolumn{6}{r}{Total recorded spend} & 1{,}832.40 \\
\bottomrule
\end{tabular}
\end{table}

Union coverage across all seventeen runs is \textbf{40/100}: every problem ever
solved by any probe was also solved by at least one of the two full passes, so
the small runs added breadth of \emph{evidence}, not breadth of \emph{coverage}.

% ============================================================
\section{Head-to-head: decomposition vs.\ the flat agent}
\label{sec:headtohead}
% ============================================================

The campaign's central comparison is \run{LUF-Anonymous} (Ultra Fleeting:
blueprint $\to$ parallel node provers $\to$ refinement) against
\run{LC-II-Irrationality} (Control~II: one continuous agent with verifier +
search, no decomposition). Both ran claude-sonnet-5
\cite{anthropic2026sonnet5}, 60-minute budgets, over the same 98 items.

\begin{table}[H]
\centering\small
\caption{Full-pass comparison. Solve time is over proved items only.}
\label{tab:headtohead}
\begin{tabular}{@{}lrr@{}}
\toprule
 & Ultra Fleeting & Control II \\
\midrule
Proved / scored          & 28 / 98 & \textbf{38 / 98} \\
Pass rate                & 28.6\% & \textbf{38.8\%} \\
Solved by this arm only  & 2 & 12 \\
Mean (median) solve time & \textbf{827\,s} (562\,s) & 1{,}395\,s (1{,}218\,s) \\
Mean attempts per scored item & 1.24 & 1.63 \\
Median tool calls / item      & 163 & 110 \\
Median LLM calls / item       & 233 & 212 \\
Notional cost                 & \$825.53 & \$878.57 \\
Notional cost per solve       & \$29.48 & \textbf{\$23.12} \\
Verifier toolchain            & v4.32.0 & v4.29.1 \\
\bottomrule
\end{tabular}
\end{table}

\subsection{Agreement and discordance}

The two arms agree on 26 solves. Control~II additionally solved 12 problems
Ultra could not (\run{fatex\_006}, \run{027}, \run{033}, \run{034},
\run{038}, \run{039}, \run{045}, \run{051}, \run{055}, \run{059}, \run{063},
\run{068}); Ultra's unique contribution was 2 (\run{fatex\_061},
\run{fatex\_067}). Treating the 14 discordant problems as Bernoulli trials, an
exact two-sided sign test gives $p = 212/2^{14} \approx 0.013$: the gap is
unlikely to be noise, despite $n=1$ passes per arm.

\subsection{Interpretation}

Three observations qualify the headline number rather than reverse it:

\begin{enumerate}[itemsep=2pt]
\item \textbf{When decomposition works, it is fast.} Ultra's median solve took
      562\,s against Control~II's 1{,}218\,s --- parallel node-filling
      genuinely buys wall-clock speed on problems whose blueprints are right.
      Its per-run internals (median 6 nodes, 4 solved, 1 forfeited;
      median 2 blueprint iterations) show most solved problems needed only
      shallow decomposition.
\item \textbf{The failure mode is the blueprint, not the nodes.} On Ultra's
      70 misses the pipeline pays a fixed planning tax before any node prover
      runs, and a wrong skeleton caps what node provers can achieve. The flat
      agent spends the same minutes probing the compiler directly
      \cite{yao2023react}, and its higher attempt count (1.63 vs.\ 1.24)
      shows the budget being recycled into fresh whole-proof attempts ---
      which FATE-X's mid-range apparently rewards.
\item \textbf{Both arms die on the same wall.} In the top quartile the two
      arms converge (2/25 each, identical problems), so the architecture
      choice matters in the middle of the difficulty range, not at the top
      (\S\ref{sec:difficulty}).
\end{enumerate}

This mirrors our earlier single-problem evidence that some FATE-X items are
``flat-solvable and decomposition-hostile'': on \run{fatex\_003}
(a finite-index coset-representative theorem), every flat arm that attempted
it succeeded --- Control~I, Control~II, Ultra's whole-pipeline pass --- while
Surround, which \emph{must} route through a skeleton, failed it
(\S\ref{sec:stronghold}).

\subsection{Difficulty gradient}
\label{sec:difficulty}

\begin{table}[H]
\centering\small
\caption{Solves by FATE-X difficulty quartile (the benchmark is ordered by
increasing difficulty \cite{fatex2025}).}
\label{tab:bands}
\begin{tabular}{@{}lrrr@{}}
\toprule
Band (item index) & Ultra Fleeting & Control II & Items scored \\
\midrule
1--25   & 11 & 12 & 24 \\
26--50  & 8  & \textbf{14} & 24 \\
51--75  & 7  & 10 & 25 \\
76--100 & 2  & 2  & 25 \\
\bottomrule
\end{tabular}
\end{table}

Control~II's entire margin comes from the middle bands (26--75: 24 vs.\ 15).
The hardest-quartile solves are the same two problems for both arms
(\run{fatex\_083}, \run{fatex\_084}), consistent with external reports that
frontier provers solve single digits on upper FATE-X
\cite{benchmarklist_fatex}. Whatever separates items 76--100 from the rest, it
is not addressed by either architecture at a 60-minute budget.

\subsection{Solves by topic area}
\label{sec:topics}

FATE-X tags every item with a topic \cite{fatex2025}. Table~\ref{tab:topics}
breaks the two full passes down by the benchmark's own finest tag (topics with
$n\ge4$ scored items shown; the remainder are in
Appendix~\ref{app:items}).

\begin{table}[H]
\centering\small
\caption{Solves by FATE-X topic tag, both full passes.}
\label{tab:topics}
\begin{tabular}{@{}lrrr@{}}
\toprule
Topic & Scored & Ultra & Control II \\
\midrule
Localization and Decomposition of Ideals & 18 & 1 & 1 \\
Galois theory & 10 & 4 & 4 \\
Ideals and Modules & 8 & 3 & 5 \\
Integral Dependence and the Nullstellensatz & 7 & 2 & 3 \\
Tensor Product and Flatness & 7 & 0 & 2 \\
Depth, Cohen--Macaulay and Gorenstein rings & 7 & 1 & 2 \\
Heights and Dimensions & 6 & 3 & 2 \\
Group Actions and Sylow theorems & 5 & 3 & 4 \\
Regular Sequences and Regular Local Rings & 4 & 3 & 3 \\
Algebraic and Arithmetic Dynamics & 4 & 1 & 1 \\
\bottomrule
\end{tabular}
\end{table}

Two structural facts emerge. First, the benchmark's largest topic ---
localization and primary decomposition, 18 items --- is nearly untouched by
both arms (1/18 each): this single topic accounts for over a quarter of all
misses and is where FATE-X's difficulty actually lives for us, more so than
the nominal difficulty rank. Second, Control~II's margin is \emph{broad},
not deep: $+2$ on Ideals and Modules, $+2$ on Tensor/Flatness, $+1$ across
four further topics --- there is no single subject the flat agent
``understands'' better, which is consistent with its advantage being
mechanical (more whole-proof attempts per hour, \S\ref{sec:headtohead})
rather than mathematical. Ultra's only topic-level win is Heights and
Dimensions ($3$ vs.\ $2$), the home of both of its unique solves
(\S\ref{sec:cases}).

% ============================================================
\section{Case studies}
\label{sec:cases}
% ============================================================

\subsection{\run{fatex\_003}: flat-solvable, decomposition-hostile}
\label{sec:fatex003}

The clearest single-problem signal in the campaign:

\begin{lstlisting}
theorem exists_leftCoset_rightCoset_representative
    (G : Type) [Group G] (H : Subgroup G) [H.FiniteIndex] :
    ∃ S : Set G, Subgroup.IsComplement S H ∧ Subgroup.IsComplement H S := by
  sorry
\end{lstlisting}

\noindent
(``A finite-index subgroup admits a common system of representatives for its
left and right cosets.'') Every flat arm that attempted it succeeded ---
Control~I in its 30-minute probe, Control~II, and Ultra's full pass --- while
Surround, the best Stronghold arm, failed it in the same campaign
(\S\ref{sec:stronghold}). The proof is a known short argument via Hall's
theorem-style matching already latent in Mathlib \cite{mathlib2020}; what it
punishes is \emph{premature structure}: a skeleton that commits to
constructing $S$ piecewise creates holes harder than the original goal.
Ultra survives it because its blueprint stage may emit a near-trivial
skeleton; Surround's planner must emit \texttt{have}-holes, and did so
wrongly. This is the concrete instance behind the ``decomposition-hostile''
label used throughout this report.

\subsection{\run{fatex\_061} and \run{fatex\_067}: where decomposition wins}

Ultra's two unique solves are topical neighbours (both tagged Heights and
Dimensions):
\begin{center}\small
\run{formallySmooth\_of\_formallySmooth\_quotient}\\
\run{isEmpty\_mvPowerSeries\_tensor\_mvPowerSeries\_algEquiv}
\end{center}
\noindent
--- formal smoothness descends along a square-zero quotient, and there is no
$k$-algebra isomorphism $k[[x]] \otimes_k k[[y]] \cong k[[x,y]]$. Both have
natural lemma decompositions --- an ideal-power calculation and a
cardinality/completeness obstruction respectively --- and both defeated the
flat agent's iterate-against-the-compiler style within budget. They are the
existence proof that the architect pipeline adds \emph{some} capability;
the finding of \S\ref{sec:headtohead} is that this capability is currently
narrow (2 problems) relative to its planning tax (10 problems).

\subsection{\run{fatex\_083} and \run{fatex\_084}: the only top-quartile solves}

Both arms' hardest-band solves are the same pair:
\begin{center}\small
\run{IsRegularLocalRing.flat\_local\_of\_regular}\\
\run{exists\_directSum\_free\_free\_of\_projective}
\end{center}
\noindent
--- regularity descends along a flat local homomorphism, and Eilenberg's
swindle (for projective $M$ there is a free $N$ with $M \oplus N$ free). Both are
textbook-named theorems whose informal proofs are certainly in the training
distribution \cite{anthropic2026sonnet5}, and both have short Mathlib-adjacent
formal skeletons. That the \emph{same} two problems fall to two very
different architectures suggests the top quartile is currently gated by
statement-level familiarity, not by search strategy --- consistent with the
contamination caveat of \S\ref{sec:validity} and with the flat difficulty
profile both arms show above item 76.

% ============================================================
\section{The control ladder}
\label{sec:controls}
% ============================================================

The controls isolate, one variable at a time, where the stack's capability
actually comes from:

\begin{description}[itemsep=2pt]
\item[Control I --- compiler only.] One continuous agent, Leak~IV verification,
  no search. \run{LC-Intuition} (30\,min, $n{=}5$ probe): 2/5, including
  \run{fatex\_003} and \run{fatex\_006}, at \$7.13 total.
\item[Control II --- compiler + search.] As Control~I plus Leak~I lemma search
  \cite{loogle}. Full pass: 38/98 (\S\ref{sec:headtohead}). Its final proofs
  lean on long \texttt{have}-chains (median 38 \texttt{have} cases per
  recorded proof), i.e.\ the flat agent \emph{internalizes} the skeleton the
  architect builds externally.
\item[Control III --- blind prover.] No tools at all: the prover emits a
  complete script from reasoning alone; a separate Leak~IV gate returns only
  pass/fail, never an error message. This isolates the value of error
  \emph{content} (Controls I/II see full compiler output; III sees one bit).
\end{description}

\subsection{Control III status and mid-campaign redesign}
\label{sec:c3}

\run{LC-III-Fair} (60\,min) is in progress: 1/3 scored at the time of writing
(\run{fatex\_002} proved on the first attempt in ${\sim}300$\,s; 95 items
pending). These early items ran the \emph{amnesiac} design: each rejected
attempt spawned a fresh session that was told only its failure count.

On 14~August the arm was redesigned (commit \run{75f6491}): all attempts now
share \emph{one continuous conversation}, resumed on every retry, with the CLI
auto-compacting history when the context window fills. The prover therefore
remembers its own past attempts --- but still never learns \emph{why} one
failed. This completes a clean three-point ladder on the memory/feedback
plane:

\begin{center}\small
\begin{tabular}{@{}lcc@{}}
\toprule
Arm & Cross-attempt memory & Error content \\
\midrule
Blind resampling (old Control III) & none & none \\
Control III (current) & full history & pass/fail bit \\
Control II & full history & full compiler output \\
\bottomrule
\end{tabular}
\end{center}

Rows recorded before and after \run{75f6491} are different arms and will be
reported separately; the pending 95 items will run the continuous design.

% ============================================================
\section{Stronghold family results}
\label{sec:stronghold}
% ============================================================

The Stronghold arms attack a flat skeleton of \texttt{have}-holes over the
Leak~II ghost-state daemon \cite{shen2026snapshot,aniva2024pantograph}. All
probes are small; treat rates as directional.

\begin{table}[H]
\centering\footnotesize
\caption{Stronghold-family probes (all claude-sonnet-5, Leak~IV gate,
v4.29.1). Calls are medians per research row.}
\label{tab:stronghold}
\begin{tabular}{@{}llrrrr@{}}
\toprule
Arm & Runs & Scored & Proved & Tool calls & LLM calls \\
\midrule
\run{have-surround} & \run{LSS-Imminent} & 4 & \textbf{3} & 95 & 174 \\
\run{have-tree} (Dark) & \run{LSD-Artist} & 1 & 0 & 71 & 138 \\
\run{finality-1} & 6 runs (\run{LFI-*}) & 9 & 1 & ${\sim}190$ & ${\sim}330$ \\
\run{stronghold-force} & \run{LSF-Patient}, \run{LFI-Undone} & 2 & 0 & 207--624 & 380--1{,}118 \\
\run{stronghold-forte} & \run{LSF-Curiosity} & 1 & 0 & 242 & 462 \\
\run{stronghold-impenetrable} & \run{LSI-Unchartered} & 0 & --- & --- & --- \\
\bottomrule
\end{tabular}
\end{table}

Three findings:

\begin{enumerate}[itemsep=2pt]
\item \textbf{Surround remains the family's best arm.} 3/4 on the opening
      slice (\run{fatex\_001}, \run{002}, \run{004}), with its one miss being
      \run{fatex\_003} --- the problem Control~I solved directly. Every
      later arm was built on top of Surround, and none has yet matched it.
\item \textbf{Complexity did not pay.} Finality-1 (timed refinement + lemma
      bank), Force, and Forte together went 1/12 while roughly doubling tool
      and LLM call counts over Surround. \run{LFI-Undone}'s single item
      consumed 624 tool calls and 1{,}118 LLM calls without a proof --- the
      machinery is exercising itself, not the problem. A capability audit
      after the Force probes found it had silently lost 11 of Finality's 13
      capabilities in porting; family arms are now diffed
      capability-by-capability against their richest sibling before
      benchmarking.
\item \textbf{The isolated-minion control (Dark) solved nothing}, consistent
      with the family's premise that minions need either a shared work-queue
      (Surround) or shared memory to make progress; but at $n=1$ this is a
      hypothesis, not a result.
\end{enumerate}

The run-scoped lemma pool (every Leak~IV-verified side lemma banked and
IDF-ranked back to minions; commits \run{5e2a250}, \run{4ebce86}) was live for
the later Stronghold probes; at current sample sizes its effect is not
separable.

% ============================================================
\section{Architect internals}
\label{sec:architect}
% ============================================================

Per-attempt research rows expose the pipeline's internal economy.

\begin{table}[H]
\centering\small
\caption{Architect-pipeline internals (medians over research rows).}
\label{tab:internals}
\begin{tabular}{@{}lrrrr@{}}
\toprule
 & \run{LUF-Anonymous} & \run{LUF-Flowing} & \run{LRV-Indomitable} \\
\midrule
Rows (attempts)        & 118 & 5 & 5 \\
Verified               & 30 & 3 & 0 \\
Nodes: total / solved / forfeited & 6 / 4 / 1 & 7 / 6 / 0 & 5 / 2 / 2 \\
Blueprint iterations   & 2 & 2 & 6 \\
Tool calls             & 163 & 116 & 178 \\
LLM calls              & 233 & 179 & 188 \\
Toolchain              & v4.32.0 & v4.32.0 & v4.32.0 \\
\bottomrule
\end{tabular}
\end{table}

Two points stand out. First, successful Ultra attempts settle their blueprint
fast (median 2 iterations) and forfeit little; \run{LRV-Indomitable}'s
failures show the opposite signature --- median 6 blueprint iterations, 2 of 5
nodes forfeited --- the pipeline churning on planning. Second,
\run{river-vintage}'s watcher team (interceptor + mechanic + consultant) did
not rescue it at 0/5: with a 20-minute budget, a 5-iteration cap, and the
older opus-4.8 driver, oversight overhead competes directly with proving time.
The vintage arm has not yet been re-benchmarked at the 60-minute budget, so
driver, budget, and watchers remain confounded in that result.

% ============================================================
\section{Cost analysis}
\label{sec:cost}
% ============================================================

Total recorded spend was \$1{,}832.40, of which the two full passes account
for 93\% (\$1{,}704). All Claude-driver figures are notional API-equivalents
reported by the CLI --- the operator pays a flat subscription, so these
numbers measure compute intensity, not cash. The one real-money run
(\run{LRV-Indomitable}, Grok driver \cite{xai2025grok41}) cost \$1.08 for
five attempts.

At the full-pass level the flat agent is cheaper per solve (\$23.12 vs.\
\$29.48) \emph{and} better, so on FATE-X there is currently no
cost-quality trade-off to argue about: Control~II dominates. The
mid-campaign cost-cap incident is worth recording: architect runs were
briefly gated by a dollar ceiling (\run{ARCHITECT\_MAX\_COST\_USD})
intended for real Grok spend; applied to the Claude driver's notional
figures it cut off healthy runs arbitrarily and has been removed for that
driver --- wall clock and iteration budgets are the real governors.

Failed attempts dominate spend: Ultra's 70 misses at roughly 60 minutes each
are ${\sim}70$ budget-hours of compute for zero score. Cheap early rejection
of unpromising problems (a cost estimator over proof-cost history is in
development) is the highest-leverage efficiency work identified by this
campaign.

% ============================================================
\section{Operational incidents and data integrity}
\label{sec:incidents}
% ============================================================

A benchmark is only as good as its bookkeeping. Three incidents from this
campaign shaped the protocol and are documented here for reproducibility:

\begin{enumerate}[itemsep=2pt]
\item \textbf{Resume-path data loss (found 8~Aug, fixed commit
      \run{f20f130}).} The console's ``from here'' requeue reset every
      non-running item at/after the target to \emph{pending} --- including
      already-scored ones. Because requeue paths never null the stored
      proof/metrics columns, the 19 wiped items of \run{LUF-Anonymous}
      (\run{fatex\_071}--\run{089}: 2 proved, 17 unsolved) were recovered by
      restoring status from the preserved columns, after a full-run JSON
      backup. The fix restricts tail-requeue to pending/skipped items.
      Table~\ref{tab:ledger} reflects post-recovery truth.
\item \textbf{Certificate wipe on failed re-verification (fixed).} A failed
      re-verify of an already-proven problem could clear its stored
      certificate before staging the new result, destroying a verified proof
      on infrastructure flake. Re-verification now stages before it clears.
\item \textbf{Verifier schema mismatch.} Leak~IV's verify endpoint takes
      \run{\{script\}} while Leak~XIV takes \run{\{code,\dots\}}; mis-wiring
      surfaces as what looks like a bridge bug and, worse, as a
      wrong-toolchain gate. Run manifests now log the armed verifier and
      toolchain line per run, which is how the v4.29.1/v4.32.0 split in
      \S\ref{sec:verifiers} is known rather than assumed.
\end{enumerate}

The triage rules of \S\ref{sec:protocol} (release vs.\ retry vs.\ skip) exist
because of a fourth, older failure class: infrastructure flake silently
consuming retry budget. Releasing a claim now decrements the attempt counter,
so no quota outage can burn a problem's attempts.

% ============================================================
\section{Threats to validity}
\label{sec:validity}
% ============================================================

\begin{enumerate}[itemsep=2pt]
\item \textbf{Single passes.} Each full-pass cell is $n{=}1$ per problem; the
      sign test in \S\ref{sec:headtohead} conditions on discordant pairs but
      cannot correct for run-level luck (service latency, provider load).
      Stronghold probes are $n{\le}4$ and merely directional.
\item \textbf{Sanitisation.} LUF benchmark run had a pre-sanitsation round to ensure 
      theorems are in the shape LUF expects (everything compacted into one theorem signature,
      statements all elaborate), but this process was omitted for LC-II benchmark run.
      It is worth mentioning LC-II's toolchain was one minor version ahead of
      the FATE-X pin. 
\item \textbf{Toolchain confound.} Ultra verified against Mathlib
      \texttt{v4.32.0}, Control~II against \texttt{v4.29.1}
      (\S\ref{sec:verifiers}); FATE-X itself pins \texttt{v4.28.0}
      \cite{fatex2025}. Lemma renames and automation changes between Mathlib
      versions can shift item difficulty in either direction. The two
      excluded items are excluded from \emph{both} arms, so the headline
      comparison is internally consistent, but absolute rates are not
      directly comparable with the FATE-X leaderboard without further
      restrictions for fairness.
      \cite{benchmarklist_fatex}.
\item \textbf{Contamination.} FATE-X sources include textbooks and exam
      archives that plausibly appear in model training data, stated by
      the authors of the benchmark in their paper.
      \cite{anthropic2026sonnet5}; formal statements are recent, informal
      ancestors may not be. This affects absolute rates, though it should
      affect both arms symmetrically.
\item \textbf{Notional costs.} Claude-driver dollar figures are
      API-equivalent estimates, not actual spend; for the reported benchmarks,
      we used Claude's cheaper Max Plan (20x usage provided comfortable headroom for 
      parallel benchmarking).
\end{enumerate}

% ============================================================
\section{Conclusions and next steps}
\label{sec:nextsteps}
% ============================================================

On FATE-X, at current budgets, our decomposition pipelines do not beat a
well-tooled flat agent: Control~II leads Ultra Fleeting 38.8\% to 28.6\% on
identical budgets ($p\approx0.013$), at lower cost per solve, and everything
any strategy has ever solved (40/100 union) is contained in what the two
full passes solved. Decomposition earns its keep in speed (562\,s vs.\
1{,}218\,s median per solve) and on exactly two problems the flat agent
missed. The Stronghold family's added machinery beyond Surround has so far
subtracted capability, and both leading arms hit the same wall in the
benchmark's top quartile.

The campaign's planned follow-ups, in priority order:

\begin{enumerate}[itemsep=2pt]
\item \textbf{Finish Control III} (95 pending items, continuous-session
      design, commit \run{75f6491}), completing the memory/feedback ladder of
      \S\ref{sec:c3}: how far does memory alone --- one bit of feedback per
      attempt --- close the gap to Control~II's 38.8\%?
\item \textbf{A full Surround pass at 60 minutes.} The family's best arm has
      only ever been probed ($n{=}4$); its 3/4 with the \run{fatex\_003}
      failure mode is the family's most interesting open signal.
\item \textbf{Toolchain-matched re-run.} Re-verify one full pass on the other
      verifier group to bound the v4.29.1/v4.32.0 confound directly.
\item \textbf{Blueprint-quality gate.} Ultra's losses concentrate in planning
      (\S\ref{sec:architect}); a cheap blueprint sanity check before
      committing node budget (cf.\ subgoal-quality filtering in
      \cite{deepseekproverv2}) is the pipeline's most promising fix.
\item \textbf{Top-quartile study.} Items 76--100 resist both architectures;
      qualitative failure analysis on \run{fatex\_085}--\run{100} should
      precede any further architecture work aimed at them.
\end{enumerate}

\newpage
\printbibliography

\newpage
\appendix
% ============================================================
\section{Per-problem results}
\label{app:items}
% ============================================================

Full-pass outcome for every FATE-X item, in the benchmark's difficulty order
\cite{fatex2025}. UF = \run{LUF-Anonymous} (Ultra Fleeting), C2 =
\run{LC-II-Irrationality} (Control~II); \cmark{} = proved, -- = unsolved
within budget. Items 9 and 41 do not elaborate under the campaign toolchains
(\S\ref{sec:benchmark}) and were excluded. Theorem names are FATE-X's own;
topics are the benchmark's finest tag.

\begingroup\footnotesize
\setlength{\tabcolsep}{4pt}
\begin{longtable}{@{}rllcc@{}}
\caption{Per-problem outcomes for the two full passes.}
\label{tab:allitems}\\
\toprule
\# & Theorem & Topic & UF & C2 \\
\midrule
\endfirsthead
\toprule
\# & Theorem & Topic & UF & C2 \\
\midrule
\endhead
\bottomrule
\endfoot
1 & \run{isPrincipalIdealRing\_of\_associated\dots} & PID, ED and UFD & \cmark & \cmark \\
2 & \run{card\_minimal\_normal\_subgroup\_le\_2} & Subgroups and Quotient g\dots & \cmark & \cmark \\
3 & \run{exists\_leftCoset\_rightCoset\_repres\dots} & Subgroups and Quotient g\dots & \cmark & \cmark \\
4 & \run{add\_one\_eq\_two\_pow\_of\_sylow\_subgro\dots} & Group Actions and Sylow \dots & \cmark & \cmark \\
5 & \run{maximal\_abelian\_normal\_subgroup\_of\dots} & p-Groups, Nilpotent Grou\dots & \cmark & \cmark \\
6 & \run{not\_isSimpleGroup\_of\_card\_eq\_396} & Group Actions and Sylow \dots & -- & \cmark \\
7 & \run{not\_isSimpleGroup\_of\_card\_eq\_1785} & Group Actions and Sylow \dots & \cmark & \cmark \\
8 & \run{quaternionAlgebra\_isomorphic\_to\_ma\dots} & Group Actions and Sylow \dots & \cmark & \cmark \\
9 & \run{sylow\_subgroup\_not\_normal\_of\_maxim\dots} & Polynomials & \multicolumn{2}{c}{excl.} \\
10 & \run{isPrincipalIdealRing\_quot\_X\_pow\_tw\dots} & Algebraic Number Theory \dots & -- & -- \\
11 & \run{not\_isomorphic\_euclideanDomain} & Basic Definitions and Ex\dots & -- & -- \\
12 & \run{isPrincipalIdealRing\_of\_quadratic\_\dots} & Integral Dependence and \dots & -- & -- \\
13 & \run{sq\_eq\_self\_of\_not\_unit} & Galois theory & -- & -- \\
14 & \run{isomorphic\_real\_of\_fractionRing\_is\dots} & Galois theory & \cmark & \cmark \\
15 & \run{exists\_swap\_stepwiseQuotient} & Galois theory & \cmark & \cmark \\
16 & \run{normalClosure\_isPExtension\_of\_isPE\dots} & Explicit Computations & -- & -- \\
17 & \run{countable\_index\_of\_maximal\_subfiel\dots} & Galois theory & -- & -- \\
18 & \run{isEmpty\_embedding\_intermediateFiel\dots} & Galois theory & -- & -- \\
19 & \run{galoisGroup\_iso\_quaternion\_group} & Algebraic Number Theory \dots & -- & -- \\
20 & \run{generated\_single\_elem\_of\_degree\_le\_p} & Galois theory & -- & -- \\
21 & \run{intermediateField\_rank\_eq\_ringChar} & Galois theory & \cmark & \cmark \\
22 & \run{finite\_algebraic\_integers\_of\_finit\dots} & Galois theory & \cmark & \cmark \\
23 & \run{ratio\_tendsto\_one\_of\_irreducible} & Galois theory & -- & -- \\
24 & \run{galoisGroup\_iso\_of\_distinct\_primes} & Galois theory & -- & -- \\
25 & \run{fixedField\_eq\_algebra\_adjoin} & Integral Dependence and \dots & -- & -- \\
26 & \run{card\_one\_or\_two\_of\_finite\_closed\_s\dots} & Integral Dependence and \dots & -- & -- \\
27 & \run{isGalois\_and\_rank\_eq\_of\_isPrimitiv\dots} & Integral Dependence and \dots & -- & \cmark \\
28 & \run{infinite\_conj\_of\_ne\_1\_absoluteGalo\dots} & Localization and Decompo\dots & -- & -- \\
29 & \run{isClosed\_zpowers\_iff\_isOfFinOrder} & Localization and Decompo\dots & \cmark & \cmark \\
30 & \run{integrallyClosedIn\_of\_complement\_m\dots} & Integral Dependence and \dots & \cmark & \cmark \\
31 & \run{UFD\_of\_ge\_5} & Localization and Decompo\dots & -- & -- \\
32 & \run{UFD\_of\_adicCompletion\_UFD} & Localization and Decompo\dots & -- & -- \\
33 & \run{isNoetherianRing\_of\_fg\_of\_isNoethe\dots} & PID, ED and UFD & -- & \cmark \\
34 & \run{powerSeries\_not\_integrallyClosed\_o\dots} & Tensor Product and Flatn\dots & -- & \cmark \\
35 & \run{noetherian\_of\_prime\_ideals\_fg} & Ideals and Modules & \cmark & \cmark \\
36 & \run{associatedPrimes\_hom\_eq\_support\_in\dots} & Integral Dependence and \dots & \cmark & \cmark \\
37 & \run{ufd\_quotDetSubOne} & Integral Dependence and \dots & -- & -- \\
38 & \run{free\_dualNumber\_and\_rank\_eq\_2} & Tensor Product and Flatn\dots & -- & \cmark \\
39 & \run{finite\_primes\_lies\_over\_of\_finite\_\dots} & Heights and Dimensions & -- & \cmark \\
40 & \run{free\_of\_rank\_iff} & Regular Sequences and Re\dots & \cmark & \cmark \\
41 & \run{isNoetherianRing\_and\_krullDim\_eq\_top} & Tensor Product and Flatn\dots & \multicolumn{2}{c}{excl.} \\
42 & \run{nonEmpty\_ringEquiv\_of\_sub\_in\_cube} & Tensor Product and Flatn\dots & -- & -- \\
43 & \run{IsRegularLocalRing.frobenius\_flat} & Ideals and Modules & -- & -- \\
44 & \run{quot\_x2\_sub\_y2\_y2\_sub\_z2\_xy\_yz\_zx\_\dots} & Completions and Hensel's\dots & -- & -- \\
45 & \run{free\_of\_free\_resolution} & Completions and Hensel's\dots & -- & \cmark \\
46 & \run{module\_flat\_iff} & Ideals and Modules & \cmark & \cmark \\
47 & \run{isEmpty\_isomorphism\_UFD\_of\_quotient} & Depth, Cohen-Macaulay Ri\dots & -- & -- \\
48 & \run{isAbsolutelyFlat\_iff\_principal\_ide\dots} & Smoothness and the Modul\dots & \cmark & \cmark \\
49 & \run{principal\_ideal\_idempotent\_iff\_fg\_\dots} & Ideals and Modules & \cmark & \cmark \\
50 & \run{isNoetherianRing\_of\_isLocalRing\_of\dots} & Tensor Product and Flatn\dots & -- & -- \\
51 & \run{adicCompletion\_faithfullyFlat\_iff} & Polynomials & -- & \cmark \\
52 & \run{generate\_unit\_ideal\_of\_quotient} & Regular Sequences and Re\dots & \cmark & \cmark \\
53 & \run{isGorensteinRing\_quot\_x2\_sub\_y2\_y2\dots} & Tensor Product and Flatn\dots & -- & -- \\
54 & \run{not\_zero\_divisor\_of\_hausdorff\_of\_d\dots} & Basic Definitions and Ex\dots & \cmark & \cmark \\
55 & \run{stablyFree\_iff\_free\_of\_not\_fg} & Ideals and Modules & -- & \cmark \\
56 & \run{projective\_of\_faithfullyFlat\_base\_\dots} & Smoothness and the Modul\dots & -- & -- \\
57 & \run{completely\_integrally\_closed\_polyn\dots} & Dedekind Domains and DVRs & \cmark & \cmark \\
58 & \run{flat\_of\_flat\_over\_quotient} & Homological Methods & -- & -- \\
59 & \run{exists\_unique\_valuation\_eq} & Polynomials & -- & \cmark \\
60 & \run{invertible\_iff\_codimension\_one} & Depth, Cohen-Macaulay Ri\dots & -- & -- \\
61 & \run{formallySmooth\_of\_formallySmooth\_q\dots} & Heights and Dimensions & \cmark & -- \\
62 & \run{adicCompletion\_equiv\_of\_smooth} & Tensor Product and Flatn\dots & -- & -- \\
63 & \run{surjection\_of\_formally\_unramified} & Ideals and Modules & -- & \cmark \\
64 & \run{homogeneous\_coordinate\_ring\_not\_is\dots} & Depth, Cohen-Macaulay Ri\dots & -- & -- \\
65 & \run{Polynomial.isGorensteinRing} & Homological Methods & -- & -- \\
66 & \run{exists\_eq\_ofList\_and\_isRegular\_of\_\dots} & Depth, Cohen-Macaulay Ri\dots & -- & -- \\
67 & \run{isEmpty\_mvPowerSeries\_tensor\_mvPow\dots} & Heights and Dimensions & \cmark & -- \\
68 & \run{localization\_jacobson\_of\_one\_lt\_ri\dots} & Depth, Cohen-Macaulay Ri\dots & -- & \cmark \\
69 & \run{IsRegularLocalRing.regularAtPrime} & Depth, Cohen-Macaulay Ri\dots & \cmark & \cmark \\
70 & \run{exists\_minimalPrime\_map\_zero} & Regular Sequences and Re\dots & \cmark & \cmark \\
71 & \run{fixedPoints\_isCohenMacaulayRing} & Regular Sequences and Re\dots & -- & -- \\
72 & \run{isCohenMacaulay\_extendScalars\_over\dots} & Completions and Hensel's\dots & -- & -- \\
73 & \run{mvPolynomial\_quotient\_isCohenMacau\dots} & Tensor Product and Flatn\dots & -- & -- \\
74 & \run{IsRegularLocalRing.gorensteinAtReg\dots} & Depth, Cohen-Macaulay Ri\dots & -- & -- \\
75 & \run{gradedAlgebra\_isCohenMacaulay\_iff\_\dots} & Ideals and Modules & -- & -- \\
76 & \run{IsCatenary.of\_noetherian\_ufd\_of\_di\dots} & Heights and Dimensions & -- & -- \\
77 & \run{infinite\_intermediate\_primes} & Ideals and Modules & -- & -- \\
78 & \run{IsCohenMacaulayLocalRing.isGorenst\dots} & Localization and Decompo\dots & -- & -- \\
79 & \run{IsLocalRing.not\_isCompleteIntersec\dots} & Group Actions and Sylow \dots & -- & -- \\
80 & \run{isCohenMacaulayRing\_of\_dimension\_t\dots} & Polynomials & -- & -- \\
81 & \run{generated\_by\_regular\_sequence\_iff} & Homological Methods & -- & -- \\
82 & \run{subring\_iso\_mvPowerSeries\_over\_DVR} & Heights and Dimensions & -- & -- \\
83 & \run{IsRegularLocalRing.flat\_local\_of\_r\dots} & Heights and Dimensions & \cmark & \cmark \\
84 & \run{exists\_directSum\_free\_free\_of\_proj\dots} & Algebraic and Arithmetic\dots & \cmark & \cmark \\
85 & \run{exist\_euclideanDomain\_not\_norm\_nat} & Algebraic and Arithmetic\dots & -- & -- \\
86 & \run{dimension\_convex} & Algebraic and Arithmetic\dots & -- & -- \\
87 & \run{exists\_polynomial\_ringEquiv\_isEmpt\dots} & Algebraic and Arithmetic\dots & -- & -- \\
88 & \run{quotient\_not\_UFD} & Localization and Decompo\dots & -- & -- \\
89 & \run{not\_isSimpleGroup\_of\_card\_eq\_336} & Localization and Decompo\dots & -- & -- \\
90 & \run{not\_finiteType\_inf\_algebraMap\_range} & Localization and Decompo\dots & -- & -- \\
91 & \run{pic\_three\_lines} & Localization and Decompo\dots & -- & -- \\
92 & \run{dimension\_sequences\_of\_one\_dimensi\dots} & Localization and Decompo\dots & -- & -- \\
93 & \run{exists\_integral\_finiteType\_not\_fin\dots} & Localization and Decompo\dots & -- & -- \\
94 & \run{zeroSet\_finite\_or\_contain\_arithmet\dots} & Localization and Decompo\dots & -- & -- \\
95 & \run{p\_map\_ne\_p} & Localization and Decompo\dots & -- & -- \\
96 & \run{ratFunc\_square\_is\_poly\_of\_orbit\_co\dots} & Localization and Decompo\dots & -- & -- \\
97 & \run{exists\_point\_not\_in\_zero\_set} & Localization and Decompo\dots & -- & -- \\
98 & \run{exists\_maximal\_ideal\_not\_in\_finite\dots} & Localization and Decompo\dots & -- & -- \\
99 & \run{equiv\_of\_aut\_equiv} & Localization and Decompo\dots & -- & -- \\
100 & \run{free\_of\_countably\_generated\_projec\dots} & Localization and Decompo\dots & -- & -- \\

\end{longtable}
\endgroup

\end{document}
